% Reviewed translation: PDF pages 207--209 / printed pages 193--195.
% Reconstructed against the original scan; environment headings remain English.
\makeatletter
\gdef\@chapapp{\chaptername}
\gdef\thechapter{\arabic{chapter}}
\gdef\theHchapter{\arabic{chapter}}
\gdef\Hy@chapapp{chapter}
\makeatother
\ctexset{
  chapter/name = {第,章},
  chapter/number = \arabic{chapter},
  chapter/numbering = true
}
\setcounter{chapter}{2}
\chapter{求解 Stokes 问题的不可压混合有限元方法}
\label{chap:incompressible-mixed-stokes}
\counterwithin{equation}{section}
\renewcommand{\thesection}{\arabic{section}}
\renewcommand{\thesubsection}{\thesection.\arabic{subsection}}
\renewcommand{\theequation}{\thesection.\arabic{equation}}
\renewcommand{\thestatement}{\thesection.\arabic{statement}}

\MathBookSourcePage{207}
\section{一个抽象问题的混合逼近}
\label{sec:incompressible-abstract-mixed-approximation}

本节重新研究 Chapter~\ref{chap:stokes-foundations} 的
Section~\ref{sec:abstract-variational-problem} 中讨论的抽象问题,
但把它置于一个较弱的框架中, 从而导出一个(通常)不同的混合表述.
由这一表述得到的混合逼近将产生一类重要的严格不可压方法,
用于求解 Stokes 和 Navier--Stokes 方程.

\subsection{一个混合变分问题}
\label{subsec:incompressible-mixed-variational-problem}
\setcounter{equation}{0}

我们处于 Subsection~\ref{subsec:abstract-general-result} 的情形中.
回顾 $X$ 和 $M$ 是两个 Hilbert 空间,
$a(\cdot,\cdot)$ 和 $b(\cdot,\cdot)$ 分别是
$X\times X$ 与 $X\times M$ 上的两个连续双线性形式.
问题 $(Q)$ 为:

给定 $(l,\chi)\in X'\times M'$, 求 $(u,\lambda)\in X\times M$, 使
\begin{align}
a(u,v)+b(v,\lambda)&=\langle l,v\rangle
&&\forall v\in X,\label{eq:incompressible-abstract-q-first}\\
b(u,\mu)&=\langle\chi,\mu\rangle
&&\forall\mu\in M.\label{eq:incompressible-abstract-q-second}
\end{align}
与往常一样, 置
\[
V(\chi)=\{v\in X;\ b(v,\mu)=\langle\chi,\mu\rangle
\quad\forall\mu\in M\},
\qquad
V=V(0).
\]
仍假设 $a(\cdot,\cdot)$ 和 $b(\cdot,\cdot)$ 满足以下两个条件:
存在常数 $\alpha>0$, 使
\begin{equation}\label{eq:incompressible-abstract-v-ellipticity}
a(v,v)\geq\alpha\|v\|_X^2
\qquad\forall v\in V;
\end{equation}
存在常数 $\beta>0$, 使
\begin{equation}\label{eq:incompressible-abstract-inf-sup}
\sup_{v\in X}\frac{b(v,\mu)}{\|v\|_X}
\geq\beta\|\mu\|_M
\qquad\forall\mu\in M.
\end{equation}
这两个假设保证问题 $(Q)$ 及其相应的问题 $(P)$:

求 $u\in V(\chi)$, 使
\begin{equation}\label{eq:incompressible-abstract-p}
a(u,v)=\langle l,v\rangle
\qquad\forall v\in V,
\end{equation}
都是适定的.

\MathBookSourcePage{208}
现在构造问题 $(Q)$ 的一个较弱表述, 它更适于我们所考虑的逼近.
引入两个自反 Banach 空间 $\widetilde X$ 和 $\widetilde M$,
其范数分别记为 $\|\cdot\|_{\widetilde X}$ 和
$\|\cdot\|_{\widetilde M}$, 并满足
\[
X\mathrel{\mathop{\subset}\limits_d}\widetilde X,
\qquad
M\mathrel{\mathop{\subset}\limits_d}\widetilde M,
\]
其中符号 $\mathrel{\mathop{\subset}\limits_d}$ 表示嵌入稠密且连续.

接着考虑两个连续双线性形式
\[
\widetilde a(\cdot,\cdot):
\widetilde X\times\widetilde X\longrightarrow\mathbb R,
\qquad
\widetilde b(\cdot,\cdot):
\widetilde X\times\widetilde M\longrightarrow\mathbb R,
\]
其范数为
\[
\|\widetilde a\|
=\sup_{u,v\in\widetilde X}
\frac{\widetilde a(u,v)}
{\|u\|_{\widetilde X}\|v\|_{\widetilde X}},
\qquad
\|\widetilde b\|
=\sup_{\substack{v\in\widetilde X\\ \mu\in\widetilde M}}
\frac{\widetilde b(v,\mu)}
{\|v\|_{\widetilde X}\|\mu\|_{\widetilde M}}.
\]
这两个双线性形式在下述意义下分别是 $a$ 和 $b$ 的延拓:
\begin{align}
\widetilde a(u,v)&=a(u,v)
&&\forall u,v\in X,\label{eq:incompressible-extended-a}\\
\widetilde b(v,\mu)&=b(v,\mu)
&&\forall v\in X,\quad\forall\mu\in M.
\label{eq:incompressible-extended-b}
\end{align}
此外, 假设 \eqref{eq:incompressible-abstract-q-first} 的右端项
$l$ 属于 $\widetilde X$ 的对偶空间 $\widetilde X'$,
并用 $\langle\cdot,\cdot\rangle$ 表示
$\widetilde X$ 与 $\widetilde X'$ 之间的对偶配对.
于是引入问题 $(\widetilde Q)$:

求 $(\widetilde u,\widetilde\lambda)
\in\widetilde X\times\widetilde M$, 使
\begin{align}
\widetilde a(\widetilde u,v)
+\widetilde b(v,\widetilde\lambda)
&=\langle l,v\rangle
&&\forall v\in\widetilde X,
\label{eq:incompressible-weak-q-first}\\
\widetilde b(\widetilde u,\mu)
&=\langle\chi,\mu\rangle
&&\forall\mu\in\widetilde M.
\label{eq:incompressible-weak-q-second}
\end{align}
对每个 $\chi\in M'$, 定义仿射集
\begin{equation}\label{eq:incompressible-weak-affine-set}
\widetilde V(\chi)
=\{v\in\widetilde X;\
\widetilde b(v,\mu)=\langle\chi,\mu\rangle
\quad\forall\mu\in\widetilde M\},
\end{equation}
以及 $\widetilde X$ 的闭子空间
\begin{equation}\label{eq:incompressible-weak-kernel}
\widetilde V=\widetilde V(0)
=\{v\in\widetilde X;\
\widetilde b(v,\mu)=0\quad\forall\mu\in\widetilde M\}.
\end{equation}
等式 \eqref{eq:incompressible-extended-b} 蕴含
\begin{equation}\label{eq:incompressible-kernel-inclusions}
V\subset\widetilde V,
\qquad
V(\chi)\subset\widetilde V(\chi).
\end{equation}
与问题 $(\widetilde Q)$ 相联系, 定义问题 $(\widetilde P)$:

求 $\widetilde u\in\widetilde V(\chi)$, 使
\begin{equation}\label{eq:incompressible-weak-p}
\widetilde a(\widetilde u,v)=\langle l,v\rangle
\qquad\forall v\in\widetilde V.
\end{equation}

\MathBookSourcePage{209}
为了便于分析问题 $(\widetilde P)$ 和 $(\widetilde Q)$,
对 $\widetilde a(\cdot,\cdot)$ 作如下假设:

形式 $\widetilde a(\cdot,\cdot)$ 是 $\widetilde V$-椭圆的,
即存在常数 $\widetilde\alpha>0$, 使
\begin{equation}\label{eq:incompressible-weak-v-ellipticity}
\widetilde a(v,v)
\geq\widetilde\alpha\|v\|_{\widetilde X}^2
\qquad\forall v\in\widetilde V.
\end{equation}
一方面, \eqref{eq:incompressible-weak-v-ellipticity}
并非 \eqref{eq:incompressible-abstract-v-ellipticity} 的推论,
因为 $\widetilde V$ 通常比 $V$ 大.
另一方面, 除了由 \eqref{eq:incompressible-abstract-inf-sup}
和 \eqref{eq:incompressible-extended-b} 推出的条件外,
并未对 $\widetilde b(\cdot,\cdot)$ 假设任何 inf-sup 条件:
\begin{equation}\label{eq:incompressible-restricted-weak-inf-sup}
\sup_{v\in\widetilde X}
\frac{\widetilde b(v,\mu)}{\|v\|_{\widetilde X}}
\geq\frac1C\sup_{v\in X}
\frac{\widetilde b(v,\mu)}{\|v\|_X}
\geq\frac\beta C\|\mu\|_M
\qquad\forall\mu\in M,
\end{equation}
其中 $C$ 是嵌入 $X\subset\widetilde X$ 的连续性常数.
严格地说, 这不足以保证问题 $(\widetilde Q)$ 适定.
下一个 Theorem 处理这一困难.

\begin{Theorem}\label{thm:incompressible-weak-mixed-equivalence}
设 $(u,\lambda)$ 是问题 $(Q)$ 的解, 且
$\widetilde a$ 满足
\eqref{eq:incompressible-weak-v-ellipticity}.

\textup{1)} 问题 $(\widetilde P)$ 在 $\widetilde V(\chi)$ 中恰有一个解
$\widetilde u$. 此外, 若 $\widetilde u$ 也属于 $V(\chi)$,
或 $V$ 在 $\widetilde V$ 中稠密, 则 $\widetilde u=u$.

\textup{2)} 若进一步有 $\lambda\in\widetilde M$,
则 $(u,\lambda)$ 是问题 $(\widetilde Q)$ 的唯一解.
\end{Theorem}

\begin{Proof}
首先回顾, Section~\ref{sec:abstract-variational-problem} 的理论也适用于
自反 Banach 空间(参见
Remark~\ref{rem:abstract-reflexive-banach-extension}).

\textup{1)} $b(\cdot,\cdot)$ 的 inf-sup 条件
\eqref{eq:incompressible-abstract-inf-sup}
蕴含 $V(\chi)$ 非空, 因而 $\widetilde V(\chi)$ 非空.
于是 $\widetilde a$ 的椭圆性蕴含问题 $(\widetilde P)$
在 $\widetilde V(\chi)$ 中有且仅有一个解 $\widetilde u$.
若 $\widetilde u\in V(\chi)$, 则由
\eqref{eq:incompressible-extended-a} 可知,
$\widetilde u$ 是问题 $(P)$ 的解; 由于 $(P)$ 恰有一个解,
故 $\widetilde u=u$.
否则, 假设 $V$ 在 $\widetilde V$ 中稠密.
于是 \eqref{eq:incompressible-abstract-p}
和 \eqref{eq:incompressible-extended-a} 蕴含
$u$ 是问题 $(\widetilde P)$ 的解, 从而 $u=\widetilde u$.

\textup{2)} 再假设 $\lambda\in\widetilde M$.
由 \eqref{eq:incompressible-extended-a}
和 \eqref{eq:incompressible-extended-b},
\eqref{eq:incompressible-abstract-q-first} 变为
\[
\widetilde a(u,v)+\widetilde b(v,\lambda)
=\langle l,v\rangle
\qquad\forall v\in X.
\]
由于 $X$ 在 $\widetilde X$ 中稠密,
这表明 $(u,\lambda)$ 是问题 $(\widetilde Q)$ 的解.
最后证明它是 $(\widetilde Q)$ 的唯一解.
第一分量 $u$ 显然唯一. 再设
\[
\widetilde b(v,\lambda)=0
\qquad\forall v\in\widetilde X.
\]
由 \eqref{eq:incompressible-restricted-weak-inf-sup}
可得 $\lambda=0$.
\end{Proof}

\begin{Remark}\label{rem:incompressible-lambda-regularity}
假设 $\lambda\in\widetilde M$ 实际上是一个正则性条件.
\end{Remark}

在后续各节中, 我们将研究问题 $(\widetilde Q)$ 的逼近比问题 $(Q)$
的逼近更简单的应用.
